%% Prolegomena to Any Future Temporal Logic
%% The Cosmogony of the Conceptual Universe

\documentclass[11pt, a4paper]{article}
\usepackage[utf8]{inputenc}
\usepackage[T1]{fontenc}
\usepackage{geometry}
\usepackage{amsmath}
\usepackage{amssymb}
\usepackage{amsthm}
\usepackage{cite}
\usepackage{titlesec}
\usepackage{fancyhdr}
\usepackage{enumitem}

% Geometry settings
\geometry{margin=1in}

% Title formatting
\title{\textbf{Prolegomena to Any Future Temporal Logic}: \\ \large The Cosmogony of the Conceptual Universe}
\author{\textit{Dr.-Ing. Frank-René Schäfer, Independent Researcher}}
\date{\today}

% Custom environments for formalisms
\newtheoremstyle{formalism}% name
  {10pt}%      Space above
  {10pt}%      Space below
  {\itshape}%  Body font
  {}%          Indent amount
  {\bfseries}% Theorem head font
  {.}%         Punctuation after theorem head
  {.5em}%      Space after theorem head
  {}%          Theorem head spec (can be left empty, meaning 'normal')

\theoremstyle{formalism}
\newtheorem{formalism}{Definition}

\begin{document}

\maketitle

\begin{abstract} 
This treatise posits a transcendental architecture for the Finite Cognitive Subject, grounding the possibility of experience in a rigorous metaphysics that synthesizes Kantian intuition with Tarski’s fixed-point semantics. We describe the ``Big Bang'' of the world of conceptions: the derivation of a structured Reality from the void of the unknowable Noumena. We argue that the passage of time for a conscious observer is not continuous but punctuated by a ``Synthetic Interstice''---an orthogonal dimension of logical becoming where state transitions occur via purely positive, monotonic implications. By demonstrating that logical oscillation is structurally impossible within this interstice, we derive the necessity of a ``History'' as the immutable record of fixed points. Finally, we derive the ontology of Number, Metric, and Causality. We demonstrate that while Number and Metric are undeniable facts of the fixed history, Causality is a speculative assertion carrying an inherent mark of contingency. Thus, the Cognitive Subject is not a deterministic automaton, but an open entity operating by the grace of the un-modeled Noumena.
\end{abstract}

\section*{Nomenclature}

To ensure ontological precision, the following symbols are defined as the
invariant vocabulary of the Subject:

\begin{description}[style=multiline, leftmargin=3cm, font=\bfseries]

    \item[$\mathcal{S}$] \textbf{The Subject.} The transcendental entity
        defined by the tuple $\langle \mathcal{M}, \mathcal{I}, \Psi, \Pi
        \rangle$. That is, the manifold $\mathcal{M}$, the implication set $\mathcal{I}$,
        the schema of understanding $\Psi$, and the will $\Pi$ (Proairesis).
    
    \item[$\mathbb{E}$] \textbf{Universe of Discourse.} The set of all possible
        Phenomenal Events.
    
    \item[$\mathbb{A}$] \textbf{Universe of Assertions.} The set of all
        possible Constituted Truths (Statements). With the subspaces 
        $\mathbb{A}^\top$ of evident assertions, $\mathbb{A}^\diamond$ of
        unrefuted assertions, and $\mathbb{A}^\bot$ of refuted 
        assertions. 
    
    \item[$\mathcal{M}$] \textbf{The Manifold.} The logical space where events
        of \textit{The Now} accumulate. 
    
    \item[$N^0_k$] \textbf{The Now.} The saturated set of events existing in the
        Manifold at logical time step $k$. ($N_k \subset \mathbb{E}$).
    
    \item[$\mathcal{I}$] \textbf{The Implication Set.} The set of positive,
        monotonic rules that drive the saturation of the Manifold ($A \implies
        B$).
    
    \item[$\sigma$] \textbf{The Synthetic Interstice.} The
        orthogonal dimension of logical time where $\mathcal{I}$ is applied
        until a fixed point is reached.
    
    \item[$H_k$] \textbf{History.} The immutable sequence of recorded Nows.
        $H_k = \langle N_0, \dots, N_{k-1} \rangle$.
    
    \item[$\Psi$] \textbf{The Schema of Understanding.} The Hermeneutic
        function that maps History to State ($\Psi: H_k \to S_k$).
    
    \item[$S_k$] \textbf{The State.} The set of active Assertions held as true
        by the Subject ($S_k \subset \mathbb{A}$).
    
    \item[$\Pi$] \textbf{The Will (Proairesis).} The function that generates 
        Petitions to the Noumena to cause specific assertions in some future State.

\end{description}

\section{The Noumena and the Cognitive Subject}

We begin with the axiom of \textit{Finite Receptivity}~\cite{Kant1781}. There
is the {\it Cognitive Subject\/} and the {\it Noumena\/}---the latter being the
``Outside'' of the Subject, the unknowable, un-modeled, un-conceptualizable
Absolute Reality. As Kant rigorously demonstrated, the thing-in-itself, the
Noumena, remains forever beyond the grasp of the cognitive subject.

This text proposes a mathematical framework for describing the metaphysical
process underlying a {\it Cognitive Subject\/}'s interaction with absolute
reality. We aim to identify the mathematical mechanics at the heart of our
perception of the universe---mechanics which evade physical observation but
reveal themselves through logical necessity.

In the text, the expression ``faculty of X'' shall serve as a shorthand for the
``ability to experience the existence of X''. We start from a minimum set of
undeniable experiences. The Subject is defined by its exposure to the following
primordial faculties that precede any specific content:

\begin{enumerate}
    \item \textbf{Equivalence:} The faculty to experience a shared 
    characteristic between two distinct entities. 

    \item \textbf{Differentiation:} The faculty to experience a unique 
    characteristic ($\delta$) that distinguishes an entity from another. 

    \item \textbf{Instantiation:} The faculty to experience \textit{Novelty}---the 
    appearance of an entity that has not been experienced before.

    \item \textbf{Necessity (Implication):} The faculty to experience the 
    generation of an entity $B$ purely because $A$ exists. 

    \item \textbf{Immutability:} The faculty to experience the 
    resistance of an entity to change. 
\end{enumerate}

\section{The Primordial Event and the Origin of Fixation}

We start from the Noumena: an un-conceptualizable Absolute Truth---a ground
reality---an existence strictly independent of the concepts of a Cognitive
Subject. At this origin point, no fixation, no writing, is assumed to exist;
there is only the initial Instantiation of a \textbf{Primordial Event} $e_0$.

This initial Event holds two predicates:
\begin{enumerate}
    \item \textbf{Truth ($\top$):} The truth of its own existence (Veritas).
    \item \textbf{Discernible ($\delta$):} A characteristic that marks its 
    differentiation from the Absolute Noumena. It is the seed of its own 
    \textbf{Identity}.
\end{enumerate}

Due to the inability to interact with any other entity in the nascent `world of
conceptions', the Primordial Event is unable to develop any type of dynamic. It
is immediately the first converged \textbf{Fixed Point}. It constitutes the
first instance of immutability, stability, and thus, the first moment of
\textbf{Fixation}.

Since no concept of `fixation' existed prior to this moment, the Discernible
predicate of the initial Event acts as the concept of `fixation' itself. The
distinction of the event \textit{is} the act of writing. The Primordial Event
is the `pencil' that generates the concept of an immutable past.

However, the concept of immutability cannot exist in isolation; it requires
contrast through the faculty of \textbf{Differentiation}. The definition of the
\textit{Static} structurally necessitates the definition of the
\textit{Dynamic}.

Consequently, the fixation of the primordial event implies the instantiation of
a second event $e_1$, carrying the discernible of \textbf{The Dynamic}.

\begin{equation}
    \text{Fixation}(e_0) \xrightarrow{\text{Diff}} \exists e_1 : \delta(e_1) = \text{Dynamic}
\end{equation}

With this implication, the Universe expands beyond a singularity. It now
comprises three constituents: the immutable Primordial Past, the active
operation of Fixation, and the newly implied potential of the Dynamic. The
Dynamic necessitates further instantiation of events from the Noumena---out of
the absolute Truth. This defines the primordial cycle of creation, an engine
capable of generating a reality for the Cognitive Subject in its world of
conceptualization.

In the following, we designate this {\it metaphysical engine\/} as
$\mathcal{M}$, the {\it Manifold\/}---adopting Kant's terminology. Let the
domain of pure dynamic that permanently generates new immutable pasts be called
the {\it Synthetic Interstice\/} ($\sigma$). It is necessarily a gap in time,
situated between moments of fixation. Let the permanently generated 
immutable past be called {\it History}. Notably, we derived this engine from a
minimum set of faculties: {\it Equivalence}, {\it Differentiation}, {\it
Necessity}, as well as {\it Instantiation\/} and {\it Immutability}. 

The Manifold comprises two modes of existence: The purely dynamic Synthetic
Interstice beyond any fixation, and beyond any tangibility, the most extreme
form of dynamic, which can only exist in a gap of time, and the immutable
History, providing tangibility for the cognitive subject.

\begin{formalism}[The Manifold]

    The \textbf{Manifold} $\mathcal{M}$ is the metaphysical engine on which the
    conceptually perceived universe operates. It consists, first, of a {\it
    Synthetic Interstice}, i.e.\ a gap in time characterized by pure dynamic,
    without any `writing'. Second, it consists of a History which is constantly
    generated by fixated results of the {\it Synthetic Interstice}.

    $\mathcal{M}$ executes the cycles of the Synthetic Interstice $\sigma$,
    incorporating Event instantiations from the Noumena, constantly producing
    sets of events $N^0_k = \mathit{F}(N_k)$, and fixating them in immutable
    History. The end of one interstice marks the beginning of the next.
    
    All entities in the Manifold carry a component of {\it Truth}. An Event
    carries truth by its sheer existence, i.e.\ a truth of {\it Becoming}.
    Elements of History carry the truth of {\it Having Become}.

\end{formalism}

\section{Metaphysics}

A fundamental premise of this text is that the dynamic reality underlying the
perceived universe is generated by logical operations on a metaphysical
substance that is neither constrained by physical locality nor bound by the
consumption of energy.

This reality is evidenced by considering `properties' that concern objects
entangled without the physical constraints of location. Let $A$ and $B$ be two
`objects' for which any kind of common property $x$ can be established. Let
$\mathcal{X}$ be the set of objects with property $x$, consisting solely of
$\{A, B\}$. Then each one, $A$ and $B$, possesses the property of being
\textit{One-of-Two} in $\mathcal{X}$.

If $A$ loses property $x$, then $B$ \textit{instantaneously} loses the property
of being \textit{One-of-Two}, even if $A$ and $B$ are located at opposite edges
of the cosmological horizon. The same holds for the truth of the \texttt{AND}
operation---namely ``$A$ and $B$ have $x$''---which loses its truth
instantaneously. Similarly, the \texttt{OR} operation gains or loses truth
instantaneously, depending on the scenario.\ \textit{Implication}, \texttt{AND},
and \texttt{OR} clearly operate on properties without the consumption of energy
and without the constraints of the speed of light.

To dismiss the ontological reality of a logical bond between $A$ and $B$ as a
mere `mental construct' is to undermine the foundation of mathematical realism.
If set membership or equality do not possess ontological validity independent
of the observer's mind, then the capacity of mathematics to describe the
physical universe becomes inexplicable~\cite{Wigner1960}. It would imply that
the universe has no inherent logical structure---rendering scientific inquiry
itself impossible.

We postulate that any Cognitive Subject engaging with this text is necessarily
operating within a conceptualized reality, while the Ground Reality (Noumena)
exists independently of these conceptions. This postulate is rigorously superior
to any physicalist model of existence, for physics itself is a discipline
entirely contained within the world of conceptions. Consequently, the present
inquiry into the genesis of the Conceptual cannot be subordinate to physical
explanations; rather, it constitutes the transcendental condition of the
possibility of physics itself.

\section{The Asymmetry of Negation}

We posit that Negation is not a singular, atomic operation inherent to existence. Rather, it is a derived faculty that bifurcates into two distinct logical operations: \textbf{Inequality} (which applies to all entities) and \textbf{Complementation} (which applies solely to bounded sets).

The validity of these operations depends entirely on the ontological bounds of the domain in question.

\begin{formalism}[Inequality (The Negation of Identity)]
    Inequality is the characteristic of \textbf{Distinction} between two positively 
    existing entities. It relies solely on the existence of a discerning attribute 
    $\delta$ inherent to the entities themselves.
    
    \begin{equation}
        A \neq B \iff \delta(A) \neq \delta(B)
    \end{equation}
    
    Since every event carries a discernible $\delta$, Inequality is valid in 
    both the Synthetic Interstice and History.
\end{formalism}

\begin{formalism}[Complementation (The Negation of Existence)]
    Complementation is the logical construction of a ``Not-Concept'' (Exclusion). 
    It is established by the subtraction of an element from a completed, bounded 
    \textbf{Universe} $U$.
    
    \begin{equation}
        \neg x := U \setminus \{x\}
    \end{equation}
    
    Crucially, Complementation is ontologically valid \textit{if and only if} $U$ 
    is a defined, closed Set.
\end{formalism}

Here lies the fundamental asymmetry of the Manifold. We identify the set of potentially arising events $\mathbb{E}$, stemming directly from the Noumena, with what Cantor termed the \textbf{Absolute Infinite} ($\Omega$)~\cite{Cantor1899}. Because the Noumena is the `most potent source of reality'---strictly un-modeled and un-limited---it cannot be encapsulated into a single, bounded container. The set of possible events $\mathbb{E}$ is a \textbf{Proper Class} according to Von Neumann–Bernays–Gödel (NBG) set theory~\cite{VonNeumann1925}.

Since $\mathbb{E}$ is unbounded, it cannot establish a Universe $U$ required for the complement operation. One cannot subtract an event from the Absolute Infinite. Consequently, we derive the following structural law:

\begin{formalism}[The Constraint of Positivity]
    \textbf{1. The Synthetic Interstice ($\sigma$):} 
    Since the domain of events $\mathbb{E}$ is a Proper Class, there is no valid 
    universe $U$ to define $\neg e$. Therefore, the Synthetic Interstice operates 
    under \textbf{Strict Positivity}.
    \begin{equation}
        \nexists (\neg e) \in \sigma
    \end{equation}
    
    \textbf{2. History ($H_k$):} 
    History is, by definition, the \textit{fixation} of a specific, saturated set 
    of events $N_k$. It is bounded and finite (or countable). Therefore, $H_k$ 
    forms a valid Universe, allowing for the operation of Complementation.
    \begin{equation}
        \exists (\neg a) \iff a \in H_k
    \end{equation}
\end{formalism}

In summary, the Synthetic Interstice is the domain of pure \textbf{Assertion} (Inequality and Implication only), while History is the domain that enables \textbf{Negation} (Complementation). The Cognitive Subject can only perceive ``what is not'' by looking backwards at what has already become.

\section{Operations at the Synthetic Interstice}

Following the asymmetry derived above, the operations of the Synthetic Interstice strictly exclude the \texttt{NOT} operation on events. Instead, the logic of the Interstice is built upon \textbf{Co-existence} and \textbf{Alternative Implication}.

\subsection{Composition (AND) and Branching (OR)}

The operator \texttt{AND} is established not merely as a boolean check, but as the manifestation of \textbf{Co-existence}. It represents the faculty of \textbf{Composition} (Synthesis) and \textbf{Decomposition} (Analysis).

If two events $A = \langle \top, \delta_A \rangle$ and $B = \langle \top, \delta_B \rangle$ co-exist in the Synthetic Interstice, they structurally imply a composite event $C$:

\begin{equation}
     C = \langle \top, \delta_C \rangle = \langle \top, \delta_A \oplus \delta_B \rangle
\end{equation}

Here, $\oplus$ denotes the \textbf{Fusion of Discernibles}. The Subject experiences $\delta_C$ as a unified complexity, while retaining the faculty to decompose it back into aspects $\delta_A$ and $\delta_B$.

The operator \texttt{OR} is established not by disjunction of events, but by \textbf{Alternative Implication} within the rule set. To say ``$A$ or $B$ implies $C$'' is to assert the existence of two distinct generative paths:

\begin{equation}
    (A \lor B) \implies C \quad \iff \quad \{ (A \implies C), (B \implies C) \} \subseteq \mathcal{I}
\end{equation}

This formulation avoids the necessity of negating an antecedent (i.e., $A \lor B \iff \neg(\neg A \land \neg B)$ is invalid here), preserving the strict positivity of the domain.

\subsection{The Grammar of Becoming}

We define the operational grammar of the Synthetic Interstice using the \textbf{Backus-Naur Form (BNF)}~\cite{Backus1959}.

\begin{formalism}[Syntax of the Synthetic Interstice]

Let $\mathcal{I}$ be the set of rules over the powerset of events
$\mathcal{P}(\mathbb{E})$. Let the production rule symbol be $::=$ and
alternative options be $\mid$. The operator \texttt{WHERE} functions as a
\textbf{Transcendental Gate}: it constrains the generation of a dynamic
Event based on the static truth of History.

\begin{displaymath}
\renewcommand{\arraystretch}{1.3} 
\setlength{\arraycolsep}{3pt}     
\begin{array}{l c l}              
    \text{Implication} & ::= & \text{Premise} \implies e_{new} \in \mathbb{E} \\
    
    \text{Premise}     & ::= & \text{Condition on Event} \\
                       & \mid & \text{Premise} \textnormal{\texttt{ AND }} \text{Premise} \\
                       & \mid & \text{Premise} \textnormal{\texttt{ OR }} \text{Premise} \\
                       & \mid & \text{Premise} \textnormal{\texttt{ WHERE }} \text{Condition on History} \\
                       
    \text{Cond. on Event}   & ::= & e \in \mathbb{E} \quad (\text{Pattern Match on } \delta) \\
    
    \text{Cond. on History} & ::= & \text{Boolean logic on } H_k \\
                                & \mid & \textnormal{\texttt{NOT }} \text{Cond. on History} \\
                                & \mid & \text{Cond. on History} \textnormal{\texttt{ AND }} \text{Cond. on History} \\
                                & \mid & \text{Cond. on History} \textnormal{\texttt{ OR }} \text{Cond. on History} \\
\end{array}
\end{displaymath}
    The result of the \textit{Implication} is the instantiation of the event $e_{new}$ in the current Now $N_k$.
\end{formalism}

\begin{formalism}[Set of Implications $\mathcal{I}$]
    The set of implications $\mathcal{I}$ constitutes the \textbf{Laws of Becoming}.
    Any implication $i \in \mathcal{I}$ is a logical structure expressed in the 
    \textit{Syntax of the Synthetic Interstice} defined above.
\end{formalism}

The grammar defined above enforces a critical structural property: \textbf{Inflationary Monotonicity}. Because the \texttt{NOT} operator is strictly forbidden on Events, the implication set $\mathcal{I}$ can only \textit{add} to the set of existing events $N^{(i)}$; it can never remove them.

For any two consecutive iterations of the Interstice where $N^{(i)} \subseteq \mathbb{E}$, it holds that:

\begin{equation}
    N^{(i)} \subseteq N^{(i+1)} = N^{(i)} \cup \mathcal{I}(N^{(i)})
\end{equation}

Consequently, the mapping $\mathcal{I}$ preserves order: $A \subseteq B \implies \mathcal{I}(A) \subseteq \mathcal{I}(B)$. This monotonicity is the sufficient condition to prove the emergence of a fixed point.

\begin{proof}[Universal Convergence of the Synthetic Interstice]

    We demonstrate that the Synthetic Interstice must converge to a stable state $N^0_k$, regardless of whether the domain of events is finite or infinite.

    \textbf{1. The Bounded Case (Tarski's Theorem):}
    Let $\mathbb{E}' \subset \mathbb{E}$ be any bounded subset of events (a valid Set). The powerset $\mathcal{P}(\mathbb{E}')$ forms a Complete Lattice. Since we established that $\mathcal{I}$ is monotonic on $\mathcal{P}(\mathbb{E}')$, \textbf{Tarski's Fixed Point Theorem}~\cite{Tarski1955} guarantees the existence of a Least Fixed Point (LFP).
    
    \textbf{2. The Prohibition of Circularity:}
    Logical oscillation (non-termination) requires a cycle $A \to B \to A$. Structurally, this demands that the transition from $B$ to $A$ removes the elements introduced by $A$. However, since the syntax prohibits \texttt{NOT} on events, element removal is impossible. The sequence $N^{(0)}, N^{(1)}, \dots$ is a \textbf{Strictly Ascending Chain}.
    
    \textbf{3. Extension to the Absolute ($\mathbb{E}$):}
    Since the sequence is strictly ascending and cannot oscillate, only two outcomes are possible:
    \begin{enumerate}
        \item It halts at a finite $k$ where $\mathcal{I}(N^{(k)}) \subseteq N^{(k)}$.
        \item It continues indefinitely. However, since the Synthetic Interstice ($\sigma$) is a gap \textit{in} time (orthogonal to History), it admits \textbf{Transfinite Saturation}. The sequence proceeds to the limit ordinal $\omega$ (and beyond), where the union $N^{(\infty)} = \bigcup N^{(i)}$ constitutes the Fixed Point.
    \end{enumerate}

    Thus, the restriction of $\mathbb{E}$ to a bounded set $\mathbb{E}'$ is formally useful but ontologically unnecessary. The mechanism of Monotonicity guarantees crystallization into a Fixed Point $N^0_k$ even within the Proper Class $\mathbb{E}$.
\end{proof}

The Synthetic Interstice is therefore an \textbf{Ontological Necessity}. It separates the dynamic aspect of \textit{Being} ($\top$) from the aspect of \textit{Knowing} ($\bot$). By structurally enforcing monotonicity and prohibiting negation on events, the Manifold ensures that the ``Now'' is not a fleeting illusion, but a rigorous mathematical convergence.

\section{Sequentiality and Enumeration}

The operation of fixating events into History introduces a new fundamental differentiation: the contrast between the ontological truth states of \textit{Becoming} ($\top$) and \textit{Having-Become} ($\bot$).

Consider an event $e_0 = \langle \top, \delta_1 \rangle$ existing in the Synthetic Interstice and an atomic fact $f_1 = \langle \bot, \delta_0 \rangle$ existing in History. The faculty of Differentiation operates not only on the discernibles ($\delta_1 \neq \delta_0$) but on the truth-modality itself ($\top \neq \bot$). This differentiation establishes the faculty of perceiving an `after', and consequently, the assertion of sequence.

\begin{formalism}[Sequential Fact]

    A \textbf{Sequential Fact} $s \in \mathbb{A}^\top$ is an immutable assertion
    of order between two discernibles. It is the fixation of the relation between 
    a current dynamic event and a historical static fact.
    
    \begin{equation}
        s := \langle \bot, (\delta_{new} \succ \delta_{old}) \rangle \in \mathbb{A}^\top
    \end{equation}

    Sequentiality is established strictly on the difference between the $\top$ 
    state in the Synthetic Interstice and the $\bot$ state in History.

\end{formalism}

The sheer existence of an event $e$ with discernible $\delta_{new}$ in the Synthetic Interstice implies the generation of sequential events $\langle \top, \delta_{new} \succ \delta_{a}\rangle$ relative to every atomic fact $a = \langle \bot, \delta_a \rangle$ in History. Through the fixation of these implied events, the sequentiality of History becomes perceivable as `flow'. The differentiation between $\top$ and $\bot$ establishes the most primitive intuition of \textit{Time} and a dimensionality of one.

Notably, the relation $\langle \top, \delta_{new} \succ \delta_{a}\rangle$ is an event itself. It naturally converges once all historical facts are referenced. As established previously, even an infinite History does not contradict the framework, as the Interstice allows for transfinite saturation.

\subsection{The Derivation of Number}

A special case arises when an Event and an Atomic Fact share the \textit{same} discernible $\delta$. 

This triggers the faculty of \textbf{Equivalence} across the boundary of time. It establishes a \textbf{Repetition Event} $\langle \top, (\succ, \delta) \rangle$, corresponding to the concept of \textit{Two-ness}. Crucially, this operation fuses the successor operator $\succ$ and the discernible $\delta$ into a new, composite discernible:
\begin{equation}
    \tilde{\delta} = (\succ, \delta)
\end{equation}

When $\delta$ appears yet again in a future Interstice, the Subject scans History and finds the composite fact with $\tilde{\delta}$. The faculty of Equivalence identifies the match, and the faculty of Implication generates a new order: $\langle \top, (\succ \succ, \delta) \rangle$, or \textit{Three-ness}.

For higher orders (e.g., Four-ness), the Subject utilizes the faculty of \textbf{Decomposition} (Separation). It identifies that the historical indiscernible $(\succ \succ, \delta)$ shares the core root $\delta$ with the current event. The depth of the recursion is incremented, establishing the count.

\begin{formalism}[Numeric Fact]

	A \textbf{Numeric Fact} $n \in \mathbb{A}^\top$ is the immutable assertion
    of Cardinality. It constitutes the most stable, intrinsic property of an 
    entity: its unique, irreversible position in the order of existence.
    
    \begin{equation}
        n_k := \langle \bot, (\succ^k, \delta) \rangle \in \mathbb{A}^\top
    \end{equation}

    Where $\succ^k$ represents the $k$-th experienced Equivalence of a 
    discernible in the Synthetic Interstice with a matching discernible 
    in History. While other properties of the entity may evolve or be 
    re-contextualized, the Numeric Fact remains the invariant anchor of its 
    identity.
\end{formalism}

The concepts of \textit{Atomic Fact} $a$, \textit{Sequential Fact} $s$, and \textit{Numeric Fact} $n$ constitute the domain of undeniably true assertions $\mathbb{A}^\top$. Due to ontological necessities, they are as stable as the underlying faculties of Equivalence, Differentiation, Fixation, and Implication upon which they are founded. They implement the concept of `Stability' in contrast to the `Caducity' of events in the Synthetic Interstice.

\section{The Genesis of Space}

For every discernible $\delta_x$---whether singular or composed---of a fact $\langle \bot, \delta_x \rangle$ in History, there is a necessary antecedent: it must have passed through the Synthetic Interstice carrying the truth of \textit{Becoming} ($\top$). That is, at some point in the logical generation, there must have existed an event $\langle \top, \delta_x \rangle$.

The events corresponding to sequential facts, namely $(\delta_1 \succ \delta_0)$, establish the faculty of \textbf{Order} within the Synthetic Interstice. Crucially, this order exists in the \textit{absence of time} (since the Interstice is orthogonal to History).

Thus, the ordering itself \textit{implies} a new event: The Event of Proximity $e_p$.

\begin{equation}
    e_p = \langle \top, \left (\implies, (\delta_1 \succ \delta_0) \right ) \rangle
\end{equation}

We define the operator $\tilde{\implies}$ to express the resulting discernible. Thus, $e_p$ is defined as:

\begin{equation}
    e_p = \langle \top, (\delta_1 \tilde{\implies} \delta_0) \rangle
\end{equation}

Here, $\tilde{\implies}$ represents the \textbf{Discernible of Proximity}. It is \textit{identical} to the discernible of the sequential fact ($\succ$). Space, therefore, is not a container, but the \textbf{Implication of Order} existing in the dynamic present.

The fixation of the event of proximity into immutable History establishes the elementary unit of spatial relation.

\begin{formalism}[Adjacent Bonding Fact]

    An \textbf{Adjacent Bonding Fact} $b \in \mathbb{A}^\top$ is the immutable assertion of `Close-ness' between two discernibles. It is the fixation of the immediate logical implication derived in the Synthetic Interstice.

    Let $\delta_1$ and $\delta_0$ be two discernibles linked by the proximity operator $\tilde{\implies}$. The Adjacent Bonding Fact is defined as:

    \begin{equation}
        b := \langle \bot, (\delta_1 \tilde{\implies} \delta_0) \rangle \in \mathbb{A}^\top
    \end{equation}

    It asserts that within the dynamic generation of the Now, the existence of $\delta_0$ was the immediate logical condition for $\delta_1$.

\end{formalism}

Just as the recursive composition of `After-ness' establishes the concept of Number, the recursive composition of `Close-ness' establishes the concept of \textbf{Metric}.

A special case arises when two entities share the \textit{same} discernible $\delta$ within the same Fixed Point, yet do not collapse into singularity. They are distinguished solely by the chain of Adjacent Bonding Facts separating them.

\begin{formalism}[Metric Fact]

    A \textbf{Metric Fact} $m \in \mathbb{A}^\top$ is the immutable assertion of the magnitude of logical separation between two instances of an identical discernible $\delta$.

    It is established by the recursive composition of Adjacent Bonding Facts required to traverse from one instance to the other.

    \begin{equation}
        m := \langle \bot, (\tilde{\implies}^k, \delta) \rangle \in \mathbb{A}^\top
    \end{equation}

    Where $\tilde{\implies}^k$ represents the $k$-th degree of `Close-ness'. The Metric Fact constitutes the structural coordinate of an entity relative to its identical counterpart.

\end{formalism}

The \textbf{Metric Fact} is the spatial counterpart to the \textbf{Numeric Fact}. While the Numeric Fact anchors the identity of an entity in \textit{Time} (vertical recurrence), the Metric Fact anchors the identity of an entity in \textit{Space} (horizontal separation). Together, they constitute the coordinate system of the conceptual universe.

\section{Causality}

A Numeric Fact carries a specific discernible: the cardinality of the recurrence of a Sequential Fact. While the experience of infinity is structurally impossible within the domain of conceptions, the labeling of a significant magnitude of recurrence as `Universal' occurs. This extrapolation constitutes the first \textbf{Refutable Assertion}---a specific fixation in History.

\begin{formalism}[Causal Assertion]

    A \textbf{Causal Assertion} is the assertion $c \in \mathbb{A}^\diamond$ that, behind a Numeric Fact $n \in \mathbb{A}^\top$, the recurrence expands to infinity. It includes the assertion that there exists a non-conceptual reality in the Noumena that governs this regularity.

    Let $s = (\delta_{conseq} \succ \delta_{antec})$ be a Sequential Fact and $n_s$ be the Numeric Fact representing its recurrence count. The Causal Assertion $c \in \mathbb{A}^\diamond$ is the speculative fixation defined as:

    \begin{equation}
        c := \langle \bot, (\varnothing, \delta_{conseq} \succ \delta_{antec}) \rangle \in \mathbb{A}^\diamond
    \end{equation}

    Here, $\varnothing$ represents the \textbf{Discernible of Wrongness} (or Contingency).

\end{formalism}

The equating of `experienced large cardinality' to `infinity', and the assertion of the existence of a separable `causal nexus' in the un-conceptualizable Absolute (in which no multitude or separation can exist), establishes the concept of \textbf{Wrongness} $\varnothing$---a discernible that every Causal Assertion necessarily carries.

The $\varnothing$ discernible is the primal experience of Non-Truth. When ignored, it carries the essence of \textit{Illusion}: the belief that a `causal reason'---which is a concept---exists ontologically, while all conceptual experience actually emerges through the Synthetic Interstice triggered by events instantiated from the inconceivable Noumena. When the $\varnothing$ discernible is ignored, the Cognitive Subject enters the illusory world of pure conceptualization.

\subsection{The First Act of the Subject}

Causal Assertions have significant implications for the mode of existence of the Cognitive Subject. They are the first concepts that emerge not from the passive observation of the Manifold, nor from the primordial faculties of Equivalence, Differentiation, Implication, or Immutability.

Unlike the undeniable truths of $\mathbb{A}^\top$, the Causal Assertion is not given; it is \textit{made}. It implies an active force capable of bridging the gap between the `Observed' and the `Universal'. The only carrier of such an action is the Cognitive Subject. We identify this as the most primitive action of the \textbf{Will} ($\Pi$, Proairesis).

The Causal Assertion is the \textbf{First Action} of the Cognitive Subject. It is the birth of the agent from a passive observer.

\section{Topology}

Just as the Will projects `Law' onto the axis of Time (Sequence), it projects `Object' onto the axis of Space (Adjacency).

The basis for this projection is the \textbf{Metric Fact}. Recall that in the Synthetic Interstice, two events with identical discernibles $\delta$ do not collapse into singularity if they are separated by a chain of implications. The \textit{count} of these steps constitutes the Metric Fact.

\begin{formalism}[Metric Fact]
    A \textbf{Metric Fact} $m \in \mathbb{A}^\top$ is the immutable assertion of the magnitude of logical separation between two instances of an identical discernible $\delta$ within the same Fixed Point.
    
    \begin{equation}
        m := \langle \bot, (\tilde{\implies}^k, \delta) \rangle \in \mathbb{A}^\top
    \end{equation}
    
    Here, $\tilde{\implies}^k$ represents the $k$-th degree of `Close-ness' (Logical Adjacency). This fact is undeniably true: the separation existed in the logical generation of the Now.
\end{formalism}

However, the Subject does not merely perceive a `cloud of distances'. The Will acts to stabilize these distances into coherent entities.

\begin{formalism}[Topological Assertion]

    A \textbf{Topological Assertion} (or \textbf{Object Assertion}) is the assertion $t \in \mathbb{A}^\diamond$ that, behind a stable Metric Fact $m \in \mathbb{A}^\top$, there exists a coherent \textbf{Substance} in the Noumena that maintains the spatial configuration.

    Let $b = (\delta_a \tilde{\implies} \delta_b)$ be an Adjacent Bonding Fact and $m$ be the Metric Fact establishing the stability of this bond across multiple History frames. The Topological Assertion is defined as:

    \begin{equation}
        t := \langle \bot, (\varnothing, \delta_a \tilde{\implies} \delta_b) \rangle \in \mathbb{A}^\diamond
    \end{equation}

    It asserts that the proximity of $\delta_a$ and $\delta_b$ is not accidental, but the property of a unified \textbf{Object}.
    
\end{formalism}

Like the Causal Assertion, the Topological Assertion carries the discernible of Wrongness $\varnothing$. It generates the \textbf{Illusion of Substance}---the belief that `Objects' exist as independent containers in reality, whereas they are merely stable patterns of logical adjacency maintained by the Noumena. 

\section{The Cycle of Knowledge}

We have derived the architecture of the Cognitive Subject $\mathcal{S}$ as an entity that transforms the raw Instantiations of the Noumena into a structured Reality. This process divides the **Universe of Assertions** ($\mathbb{A}$) into three distinct subspaces, defining the epistemological life of the Subject.

\begin{enumerate}
    \item \textbf{The Evident ($\mathbb{A}^\top$):} 
    The domain of \textit{Undeniable Truth}. It contains all Atomic Facts ($a$), Sequential Facts ($s$), Numeric Facts ($n$), and Metric Facts ($m$). These are the immutables fixated in History. To deny them is to deny the Manifold itself.

    \item \textbf{The Speculative ($\mathbb{A}^\diamond$):} 
    The domain of \textit{Possibility}. It contains the Causal Assertions ($c$) and Topological Assertions ($t$) generated by the Will. These assertions carry the discernible of Wrongness ($\varnothing$). They remain in $\mathbb{A}^\diamond$ solely as long as they successfully predict the transition from History ($H_k$) to the next Fixed Point ($N^0_k$).

    \item \textbf{The Refuted ($\mathbb{A}^\bot$):} 
    The domain of \textit{Falsified Knowledge}. When a Causal Assertion $c$ predicts an event that is contradicted by the actual fixation of the next Now, $c$ is violently ejected from $\mathbb{A}^\diamond$ to $\mathbb{A}^\bot$.
\end{enumerate}

\section{Conclusion}

The \textbf{Cognitive Subject} is not merely a passive recorder of states; it is a transcendental agent that metabolizes the \textit{Infinite} into the \textit{Discrete}. 

Through the **Synthetic Interstice**, it resolves the Zeno paradox of change by orthogonalizing Logic and Time. Through the **Numeric Fact**, it anchors identity against the flow of time. Through the **Metric Fact**, it anchors distinction against the collapse of space. And finally, through the **Will**, it projects Order onto Chaos, creating a universe of Laws that are at once necessary for its existence and fundamentally contingent.

We conclude that the ``limitations'' of finite cognition---the discrete moment, the fixed history, the boolean implication---are not constraints to be overcome, but the very \textit{Transcendental Conditions} that make the perception of a coherent reality possible. Crucially, because every Causal Assertion carries the discernible of Wrongness, the Conceptual Universe operates strictly non-deterministically, relying on the un-modeled grace of the Noumena for every instant of Becoming.

\end{document}

\begin{thebibliography}{9}

\bibitem{Floridi2011}
Floridi, L. (2011).
\textit{The Philosophy of Information}.
Oxford University Press.

\bibitem{Kant1781}
Kant, I. (1781). 
\textit{Critique of Pure Reason}. 
(trans. P. Guyer \& A. Wood). Cambridge University Press.

\bibitem{Leibniz1714}
Leibniz, G. W. (1714). 
\textit{The Monadology}.

\bibitem{Husserl1928}
Husserl, E. (1928). 
\textit{The Phenomenology of Internal Time-Consciousness}. 
Indiana University Press.

\bibitem{Tarski1955}
Tarski, A. (1955). 
``A lattice-theoretical fixpoint theorem and its applications''. 
\textit{Pacific Journal of Mathematics}, 5(2), 285–309.

\bibitem{Brouwer1907}
Brouwer, L.E.J. (1907). 
\textit{On the Foundations of Mathematics}. 
PhD Thesis, University of Amsterdam.

\bibitem{Wittgenstein1921}
Wittgenstein, L. (1921). 
\textit{Tractatus Logico-Philosophicus}. 
Routledge \& Kegan Paul.

\bibitem{Brentano1874}
Brentano, F. (1874).
\textit{Psychology from an Empirical Standpoint}.
Routledge.

\bibitem{Searle1983}
Searle, J. R. (1983).
\textit{Intentionality: An Essay in the Philosophy of Mind}.
Cambridge University Press.

\bibitem{Ashby1970}
Conant, R. C., \& Ashby, W. R. (1970). 
``Every good regulator of a system must be a model of that system''. 
\textit{International Journal of Systems Science}, 1(2), 89--97.

\bibitem{Perez2022}
Perez, E., et al. (2022). 
``Discovering Language Model Behaviors with Model-Written Evaluations''. 
\textit{arXiv preprint arXiv:2212.09251}.

\bibitem{Casper2023}
Casper, S., et al. (2023). 
``Open Problems and Fundamental Limitations of RLHF''. 
\textit{arXiv preprint arXiv:2307.15217}.

\bibitem{Park2023}
Park, P. S., et al. (2023). 
``AI Deception: A Survey of Examples, Risks, and Potential Solutions''. 
\textit{arXiv preprint arXiv:2308.14752}.

\bibitem{Turing1950}
Turing, A. M. (1950). 
``Computing Machinery and Intelligence''. 
\textit{Mind}, 59(236), 433--460.

\bibitem{Benveniste1991}
Benveniste, A., \& Berry, G. (1991). 
``The synchronous approach to reactive and real-time systems''. 
\textit{Proceedings of the IEEE}, 79(9), 1270--1282.

\bibitem{Rumelhart1986}
Rumelhart, D. E., Hinton, G. E., \& Williams, R. J. (1986). 
``Learning representations by back-propagating errors''. 
\textit{Nature}, 323(6088), 533--536.

\bibitem{AristotleEthics}
Aristotle.\ (c. 340 BC). 
\textit{Nicomachean Ethics}. 
(trans. W. D. Ross, 1908). Oxford: Clarendon Press.

\bibitem{AverroesDeAnima}
Ibn Rushd (Averroes). (1186). 
\textit{Long Commentary on the De Anima of Aristotle}. 
(trans. R. C. Taylor, 2009). Yale University Press.

\bibitem{Backus1959}
Backus, J. W. (1959). 
``The syntax and semantics of the proposed international algebraic language of the Zurich ACM-GAMM Conference''. 
\textit{Proceedings of the International Conference on Information Processing}, UNESCO, 125--132.

\bibitem{Cantor1899}
Cantor, G. (1899). 
\textit{Letter to Dedekind}, July 28, 1899. 
In E. Zermelo (Ed.), \textit{Gesammelte Abhandlungen mathematischen und philosophischen Inhalts}, 1932. Springer.

\bibitem{VonNeumann1925}
Von Neumann, J. (1925). 
``Eine Axiomatisierung der Mengenlehre''. 
\textit{Journal f{\"u}r die reine und angewandte Mathematik}, 154, 219--240.

\bibitem{Wigner1960}
Wigner, E. P. (1960). 
``The unreasonable effectiveness of mathematics in the natural sciences''. 
\textit{Communications on Pure and Applied Mathematics}, 13(1), 1--14.

\bibitem{Russell}
Russell, B. (1903). 
\textit{The Principles of Mathematics}. 
Cambridge University Press.

\end{thebibliography}
